\subsection*{The affine quotient beyond the rank-129 homology action}

\paragraph{FA.1 (Proved): the actual quotient as a path-chain extension.}
Let $C=C_2\amalg C_2\amalg C_2$, generated by
$(t_0,t_1,t_2)$, and put
\[
 D=C''(C')^2,\qquad Q=C/D,\qquad |Q|=256.
\]
The Cayley graph of $Q$ with these three involutions
has $256$ vertices and $384$ unoriented edges.
Orient each edge arbitrarily, with a negative sign
for its reversed orientation. Over $\mathbb F_3$,
write $C_1$ for the edge-chain space and
$I_Q=\ker(\mathbb F_3[Q]\to\mathbb F_3)$ for the
augmentation ideal. The boundary map gives an exact
sequence of $Q$-modules
\[
 0\longrightarrow M=H_1(D,\mathbb F_3)
 \longrightarrow C_1\stackrel{\partial}{\longrightarrow}
 I_Q\longrightarrow0,
 \qquad \dim(M,C_1,I_Q)=(129,384,255).
 \tag{FA.1a}
\]
The dimension follows either from the graph or the
index-$32$ Schreier formula inside the free rank-five
group $C'$. We first use the discrete free product and its
finite covering graph, and then complete the resulting
free finite-index subgroup. Completion preserves its
finite-coefficient abelianization. Thus this graph
computes the stated homology of the profinite subgroup
$D$ with these finite coefficients.

Let $q\cdot c$ denote deck translation of an edge chain.
Inside $C_1\rtimes Q$, define
\[
 \mathcal E=\{(c,q):\partial c=q-1\},\qquad
 (c,q)(d,r)=(c+q\cdot d,qr).
 \tag{FA.1b}
\]
This is a subgroup, since
$(q-1)+q(r-1)=qr-1$. Sending $t_i$ to its one-edge
path from $1$ to $t_i$ defines a homomorphism
$C\to\mathcal E$: the relation $t_i^2=1$ holds
because the translated second edge is the negative
of the first. The image of a word is exactly its
path chain and endpoint.

On $D$ the endpoint is one and the map is the usual
homology map onto $M$, with kernel $D'D^3$. For each
$q\in Q$, choose a path from $1$ to $q$. Its chain,
plus arbitrary cycles from $M$, gives every chain
with boundary $q-1$. Thus the homomorphism is onto
and
\[
       \boxed{\ C/(D'D^3)\simeq\mathcal E,\qquad
                    |\mathcal E|=256\cdot3^{129}.\ }
 \tag{FA.1c}
\]
This is an actual finite group quotient with its
source marking, not a linear approximation to it.

Since $|Q|$ is invertible in $\mathbb F_3$, averaging
an arbitrary linear splitting of $\partial$ gives
a $Q$-equivariant splitting $s:I_Q\to C_1$.
The maps $q\mapsto(s(q-1),q)$ form a group section
of (FA.1b). Consequently $\mathcal E\simeq M\rtimes Q$,
although the path-chain model avoids choosing that
splitting and retains the exact based marking.

\paragraph{FA.2 (Proved): centrelessness and the information lost by the linear action.}
The edge module is
\[
 C_1=\bigoplus_{i=0}^2
       \operatorname{Ind}_{\langle t_i\rangle}^{Q}
                      \operatorname{sgn}.
\]
It has no $Q$-invariants: on a fixed edge, the
stabilizing involution reverses orientation, and
$-1\ne1$ in $\mathbb F_3$. Hence
\[
                              M^Q=0.
 \tag{FA.2a}
\]
The action of $Q$ on $M$ is faithful. To prove this
without a numerical faithfulness test, use the same
integral graph cycle lattice $M_{\mathbb Z}$, whose
profinite completion is $D^{\rm ab}$. For $1\ne q\in Q$, the
vertex action has trace zero. An edge fixed by $q$
is reversed, so its trace contribution is minus one.
The homology trace is therefore
\[
          \operatorname{tr}(q\mid M_{\mathbb Z}\otimes\mathbb Q)
                       =1-\#\{q\text{-fixed edges}\}\leq1,
\]
which is strictly below its dimension $129$.
If $q$ acted as the identity modulo three, its
integral action on $\mathbb Z_3$ homology would lie
in the pro-three kernel of
$\operatorname{GL}_{129}(\mathbb Z_3)\to
\operatorname{GL}_{129}(\mathbb F_3)$. But its order
is a power of two, so that action would be the
identity, contradicting the trace. This proves
faithfulness over $\mathbb F_3$.

A central element of $\mathcal E$ consequently
has trivial $Q$ coordinate, since it centralizes
$M$; its remaining coordinate belongs to $M^Q$.
Thus
\[
                         Z(\mathcal E)=1.
 \tag{FA.2b}
\]
The source action of $L$ on $C$ preserves $D'D^3$
and therefore induces actual automorphisms of
$\mathcal E$. On chains it is obtained by replacing
each oriented edge by the literal path spelling
the image of its involution. It is semilinear
with respect to the induced permutation of $Q$
and intertwines the boundary maps. Its restriction
to $M$ is the old rank-129 homology action.

For $m\in M$, however, inner conjugation in the
full path-chain group is
\[
       (c,q)\longmapsto(c+(1-q)m,q).
 \tag{FA.2c}
\]
It is the identity on cycles, whose endpoint is
one, but need not be the identity on chains.
It is the identity everywhere if and only if
$m\in M^Q$, hence if and only if $m=0$.
In any splitting of (FA.1a), this is an upper
triangular shear from the augmentation-ideal
part to the homology part. This is exactly
the information discarded by retaining only
the linear action on $M$.

\paragraph{FA.3 (Proved formulation; the new obstruction is not evaluated here).}
For a genuine gentle multiplier-one pair, RP.8
proves $\mathcal P(f)\in D$. Its image in the
new quotient is therefore the definite vector
\[
           m_f=[\mathcal P(f)]\in D/(D'D^3)=M.
 \tag{FA.3a}
\]
The full edge-chain action of its pentagon is
the shear (FA.2c) for this vector. Thus
\[
 \mathcal P(f)=1\text{ in }C/(D'D^3)
 \quad\Longleftrightarrow\quad m_f=0
 \quad\Longleftrightarrow\quad
 \mathcal P(f)\text{ acts trivially on }C_1.
 \tag{FA.3b}
\]
The old rank-129 action is always trivial on
this element and cannot decide any of these
equivalent conditions. The exact sequence
(FA.1a) gives a concrete finite target in
which the remaining vector can be computed.
No calculation in this paragraph proves
$m_f=0$ for all gentle elements, identifies
a coordinate of $m_f$ with $\Delta$, or
promotes a finite solution to full GT.

Geometrically, this extension retains based
paths between the $256$ deck sheets, in
addition to their closed cycles. The genus-two
Prym calculations describe constituents of
the cycle module $M$; they do not supply
these based paths or their extension class.
That distinction is why their linear
boundary/proper coupling can be nontrivial
while the linear pentagon remains automatic.
